\documentclass{article}
\usepackage[utf8]{inputenc}
\usepackage{amsmath,amssymb}
\usepackage[most]{tcolorbox}
\usepackage{geometry}
\geometry{margin=2.5cm}

\newcommand{\step}[1]{\par\noindent\hangindent=3.5em\hangafter=1%
  \makebox[3.5em][l]{\textbf{#1.}}}
\newcommand{\steptag}[1]{\hfill{\small\textsf{[#1]}}}

\newtcolorbox{proofbox}[1]{
  colback=gray!5, colframe=gray!60,
  fonttitle=\bfseries, title={#1},
  breakable, enhanced,
  left=4pt, right=4pt, top=4pt, bottom=4pt,
  before skip=10pt, after skip=10pt
}

\begin{document}

\begin{proofbox}{The signed-idempotent and stochastic-idempotent formulati...}
\step{1} The signed-idempotent and stochastic-idempotent formulations of classical stability are equivalent up to universal constants: Q row-stochastic with ||Q\textasciicircum{}2-Q|| <= eta gives P=theta(2Q-1) signed affine retraction with ||P-Q|| <= C eta and neg mass delta <= C eta, and conversely row-normalising p\_i\textasciicircum{}+ gives Q with ||P-Q|| <= 2 delta, ||Q\textasciicircum{}2-Q|| <= 6 delta+4 delta\textasciicircum{}2. \steptag{validated} \\[6pt]
\step{1.1} N1 (ARENA, shared by both directions). M\_n(R) under the operator norm ||A|| := ||A|$|\_\{inf-\rangle$inf\} = sup\_\{||x||\_inf<=1\} ||Ax||\_inf, induced by ||x||\_inf=max\_k|x\_k| on R\textasciicircum{}n, is a unital Banach algebra (GT-operator-norm-banach: bounded linear operators on a Banach space form a unital Banach algebra, instantiated at calA=R\textasciicircum{}n). Concretely: ||I||=1; triangle ||A+B||<=||A||+||B|| and exact submultiplicativity ||AB||<=||A|| ||B|| follow from the sup-definition (no 1+eps slack); finite-dim => complete. CLOSED FORM (EXTRACTION-LEVEL FLAG -- NOT byte-stated in any present ref; derived inline, established lem-P-properties node 1.1.2 pattern): for ||x||\_inf<=1, |(Ax)\_i|<=sum\_j|A\_ij|, with equality at x\_j=sign(A\_ij), so ||A||=max\_i sum\_j|A\_ij| (max row ell\textasciicircum{}1). In particular ||A||=max\_i ||a\_i||\_1 where a\_i is row i; for a probability/signed row a\_i, ||a\_i||\_1 is its total variation. [GT-operator-norm-banach] \steptag{validated} \\[4pt]
\step{1.2} N2 (FORWARD spectral defect). Let Q be row-stochastic (GT-hyp / def-stochastic: Q>=0 entrywise, Q1=1) with ||Q\textasciicircum{}2-Q||<=eta, eta<=eta\_0<1/4. By N1 (1.1), ||Q||=max\_i ||q\_i||\_1 = max\_i sum\_j Q\_ij = 1 (rows are probability vectors: Q>=0 and Q1=1 give row sums 1). Set S:=2Q-I. Then S1=2(Q1)-1=2*1-1=1, and ||S||<=2||Q||+||I||=2+1=3 (N1 triangle, ||I||=1). Also S\textasciicircum{}2-I=(2Q-I)\textasciicircum{}2-I=4Q\textasciicircum{}2-4Q+I-I=4(Q\textasciicircum{}2-Q), so ||S\textasciicircum{}2-I||=4||Q\textasciicircum{}2-Q||=4 eta <= 4 eta\_0 < 1 (GT-small-eta-threshold: the eta<1/4 threshold; here ||S\textasciicircum{}2-I||=4eta<1 exactly matches Kitaev's ||X\textasciicircum{}2-I||<=4delta<1 under P->Q, delta->eta). \steptag{validated} \\[4pt]
\step{1.3} N3 (FORWARD inverse-sqrt + EXPLICIT C -- HONEST FLAG 1: the explicit constant is DERIVED here, not cited; Kitaev gives only big-O). Apply GT-funcalc-taylor-bound with B=M\_n(R) (N1), f(x)=x\textasciicircum{}\{-1/2\}, x0=1, r=1, X=S\textasciicircum{}2. Since ||S\textasciicircum{}2-I||=4eta<1=r (N2), R:=(S\textasciicircum{}2)\textasciicircum{}\{-1/2\}=sum\_\{n>=0\} a\_n (S\textasciicircum{}2-I)\textasciicircum{}n converges, commutes with S\textasciicircum{}2 (and S), with a\_n=C(-1/2,n). DERIVED scalar real-analysis (one line, no external): |a\_n|=|C(-1/2,n)|=C(2n,n)/4\textasciicircum{}n>=0 (the signs alternate so |a\_n|=(-1)\textasciicircum{}n C(-1/2,n)), and sum\_\{n>=0\}|a\_n| t\textasciicircum{}n=(1-t)\textasciicircum{}\{-1/2\} for 0<=t<1. Hence R1=a\_0 1 + sum\_\{n>=1\}a\_n (S\textasciicircum{}2-I)\textasciicircum{}n 1 = 1 (a\_0=f(1)=1; (S\textasciicircum{}2-I)1=S\textasciicircum{}2 1 - 1 = S(S1)-1 = S*1-1 = 1-1 = 0 by N2 S1=1, so all n>=1 terms vanish). By the GT-funcalc-taylor-bound sum-bound, ||R-I||=||f(S\textasciicircum{}2)-f(I)|| <= sum\_\{n>=1\}|a\_n| ||S\textasciicircum{}2-I||\textasciicircum{}n = ||S\textasciicircum{}2-I|| sum\_\{n>=1\}|a\_n| ||S\textasciicircum{}2-I||\textasciicircum{}\{n-1\} <= C\_1(eta\_0) ||S\textasciicircum{}2-I||, where C\_1(eta\_0):=sum\_\{n>=1\}|a\_n| (4 eta\_0)\textasciicircum{}\{n-1\} <= (4eta\_0)\textasciicircum{}\{-1\}((1-4eta\_0)\textasciicircum{}\{-1/2\}-1) < infinity is EXPLICIT and finite (geometric-type tail, 4eta\_0<1). Thus ||R-I|| <= C\_1(eta\_0) ||S\textasciicircum{}2-I|| = 4 C\_1(eta\_0) eta. [GT-funcalc-taylor-bound; N1; N2] \steptag{validated} \\[4pt]
\step{1.4} N4 (FORWARD calculus multiplicativity -- HONEST FLAG 2, TENSION for the verifier). CLAIM: R\textasciicircum{}2=(S\textasciicircum{}2)\textasciicircum{}\{-1\}, hence sgn(S):=S R satisfies sgn(S)\textasciicircum{}2=S\textasciicircum{}2 R\textasciicircum{}2=I exactly. DERIVED as a Cauchy-product node (conservative), NOT read off Kitaev's bytes: Kitaev:535 only DISMISSES the holomorphic calculus and does not state f(X)g(X)=(fg)(X). (TENSION the verifier must byte-check: the sibling repo's proofs/lem-P-properties external GT-funcalc (0f14e166) DOES assert 'the calculus is multiplicative f(X)g(X)=(fg)(X)' as cited @ :503-532,:642; we do NOT rely on that wording, we re-derive it.) DERIVATION: write Y:=S\textasciicircum{}2-I, ||Y||=4eta<1 (N2); R=sum\_m a\_m Y\textasciicircum{}m and R\textasciicircum{}2=(sum\_m a\_m Y\textasciicircum{}m)(sum\_k a\_k Y\textasciicircum{}k). Both series converge ABSOLUTELY in operator norm: sum\_m |a\_m| ||Y||\textasciicircum{}m <= (1-||Y||)\textasciicircum{}\{-1/2\}<infinity (N3) and submultiplicativity ||Y\textasciicircum{}\{m+k\}||<=||Y\textasciicircum{}m|| ||Y\textasciicircum{}k|| (N1), all powers of the single element Y commute, so the Cauchy product is justified and R\textasciicircum{}2 = sum\_N (sum\_\{m+k=N\} a\_m a\_k) Y\textasciicircum{}N = sum\_N b\_N Y\textasciicircum{}N where b\_N are the coefficients of f(x)\textasciicircum{}2=((1+y)\textasciicircum{}\{-1/2\})\textasciicircum{}2=(1+y)\textasciicircum{}\{-1\}=sum\_N (-1)\textasciicircum{}N y\textasciicircum{}N (y=x-1). Thus R\textasciicircum{}2=(S\textasciicircum{}2)\textasciicircum{}\{-1\}=sum\_N (-1)\textasciicircum{}N Y\textasciicircum{}N=(I+Y)\textasciicircum{}\{-1\}, and (I+Y)=S\textasciicircum{}2, so S\textasciicircum{}2 R\textasciicircum{}2=I. Therefore sgn(S)\textasciicircum{}2=(SR)(SR)=S\textasciicircum{}2 R\textasciicircum{}2=I (S,R commute via the common functional calculus in S\textasciicircum{}2, N3). [GT-funcalc-taylor-bound; N1; N3] \steptag{validated} \\[4pt]
\step{1.5} N5 (FORWARD: P is an exact signed affine retraction). Define P:=theta(2Q-I)=(1/2)(I+sgn(S)) where S=2Q-I (N2), sgn(S)=S R (N4). This is exactly the GT-kitaev-spectral-propP construction with the renaming P\_Kitaev->Q, delta->eta (||Q\textasciicircum{}2-Q||<=eta<1/4, GT-hyp): Prop\_P gives \textasciitilde{}P=theta(2Q-I) with \textasciitilde{}P\textasciicircum{}2=\textasciitilde{}P EXACTLY. We re-derive idempotence directly to keep it self-contained: P\textasciicircum{}2=(1/4)(I+sgn(S))\textasciicircum{}2=(1/4)(I+2 sgn(S)+sgn(S)\textasciicircum{}2)=(1/4)(I+2 sgn(S)+I)=(1/4)(2I+2 sgn(S))=(1/2)(I+sgn(S))=P, using sgn(S)\textasciicircum{}2=I (N4). Unit-preservation: P1=(1/2)(1+sgn(S)1)=(1/2)(1+S(R1))=(1/2)(1+S*1)=(1/2)(1+S1)=(1/2)(1+1)=1, using R1=1 (N3) and S1=1 (N2). So P1=1 and P\textasciicircum{}2=P exactly; P's rows are signed measures (real entries) of total mass 1 (from P1=1), i.e. P is a signed affine retraction in the sense of def-stochastic. [GT-kitaev-spectral-propP; N2; N3; N4] \steptag{validated} \\[4pt]
\step{1.6} N6 (FORWARD closeness, EXPLICIT universal C). Compute P-Q: P=(1/2)(I+SR), and Q=(1/2)(I+S) since S=2Q-I => Q=(1/2)(I+S). So P-Q=(1/2)(I+SR)-(1/2)(I+S)=(1/2)(SR-S)=(1/2)S(R-I). By N1 submultiplicativity and N2 (||S||<=3), N3 (||R-I||<=4 C\_1(eta\_0) eta): ||P-Q||<=(1/2)||S|| ||R-I||<=(1/2)(3)(4 C\_1(eta\_0) eta)=6 C\_1(eta\_0) eta. Define the EXPLICIT universal constant C:=6 C\_1(eta\_0) (HONEST FLAG 1: C is DERIVED, C\_1(eta\_0)=sum\_\{n>=1\}|C(-1/2,n)|(4eta\_0)\textasciicircum{}\{n-1\}<infinity from N3, depending only on the fixed regime parameter eta\_0<1/4, dimension-free in n). Hence ||P-Q||<=C eta. [N1; N2; N3; N5] \steptag{validated} \\[4pt]
\step{1.7} N7 (FORWARD negative mass; closes FORWARD). The negative mass of P is delta:=max\_i sum\_j max(-P\_ij,0) (def-stochastic). Since Q>=0 entrywise (GT-hyp / def-stochastic), for each (i,j): max(-P\_ij,0)=max(Q\_ij-P\_ij-Q\_ij,0)<=max(Q\_ij-P\_ij,0)<=|Q\_ij-P\_ij|=|P\_ij-Q\_ij| (using Q\_ij>=0 so -P\_ij<=Q\_ij-P\_ij, and max(t,0)<=|t|). Summing over j: sum\_j max(-P\_ij,0) <= sum\_j |P\_ij-Q\_ij| = ||p\_i-q\_i||\_1 (row i of P-Q). By N1 (||P-Q||=max\_i ||p\_i-q\_i||\_1): delta=max\_i sum\_j max(-P\_ij,0) <= max\_i ||p\_i-q\_i||\_1 = ||P-Q|| <= C eta (N6, same universal C). CONCLUSION (FORWARD): from row-stochastic Q with ||Q\textasciicircum{}2-Q||<=eta, P=theta(2Q-1) is a signed affine retraction (P1=1, P\textasciicircum{}2=P exactly, N5) with ||P-Q||<=C eta (N6) and negative mass delta<=C eta. [N1; N6; def-stochastic] \steptag{validated} \\[4pt]
\step{1.8} N8 (CONVERSE construction + distances). Let P be a signed affine retraction (GT-hyp / def-stochastic: P1=1, P\textasciicircum{}2=P exactly, rows p\_i signed measures of total mass 1) with negative mass delta=max\_i sum\_j max(-P\_ij,0). Disjoint-support split of each row: p\_ij=p\_ij\textasciicircum{}+ - p\_ij\textasciicircum{}- with p\_ij\textasciicircum{}+=max(p\_ij,0), p\_ij\textasciicircum{}-=max(-p\_ij,0); supports \{j:p\_ij>0\} and \{j:p\_ij<0\} are disjoint. ell\textasciicircum{}1 additivity on disjoint supports: ||a p\_i\textasciicircum{}+ - b p\_i\textasciicircum{}-||\_1=a||p\_i\textasciicircum{}+||\_1+b||p\_i\textasciicircum{}-||\_1 for a,b>=0. Total mass: sum\_j p\_ij=(P1)\_i=1 (P1=1). Set a\_i:=||p\_i\textasciicircum{}-||\_1=sum\_j max(-P\_ij,0)=neg(p\_i)<=delta. Then ||p\_i\textasciicircum{}+||\_1=sum\_j p\_ij + ||p\_i\textasciicircum{}-||\_1 = 1 + a\_i (since sum\_j p\_ij\textasciicircum{}+ - sum\_j p\_ij\textasciicircum{}- = sum\_j p\_ij =1). Define q\_i:=p\_i\textasciicircum{}+/(1+a\_i) (well-defined: 1+a\_i>=1>0). Then q\_i>=0 entrywise and ||q\_i||\_1=||p\_i\textasciicircum{}+||\_1/(1+a\_i)=(1+a\_i)/(1+a\_i)=1, so Q (rows q\_i) is row-stochastic (Q>=0, Q1=1, def-stochastic), and by N1 ||Q||=max\_i ||q\_i||\_1=1. DISTANCES: p\_i-q\_i=p\_i\textasciicircum{}+ - p\_i\textasciicircum{}- - p\_i\textasciicircum{}+/(1+a\_i)=p\_i\textasciicircum{}+ (1-1/(1+a\_i)) - p\_i\textasciicircum{}- = p\_i\textasciicircum{}+ (a\_i/(1+a\_i)) - p\_i\textasciicircum{}-, disjoint supports, so ||p\_i-q\_i||\_1=(a\_i/(1+a\_i))||p\_i\textasciicircum{}+||\_1 + ||p\_i\textasciicircum{}-||\_1=(a\_i/(1+a\_i))(1+a\_i)+a\_i=a\_i+a\_i=2 a\_i. Hence ||P-Q||=max\_i ||p\_i-q\_i||\_1=max\_i 2 a\_i<=2 delta (N1). Also ||p\_i||\_1=||p\_i\textasciicircum{}+||\_1+||p\_i\textasciicircum{}-||\_1=(1+a\_i)+a\_i=1+2 a\_i, so ||P||=max\_i ||p\_i||\_1=max\_i(1+2 a\_i)=1+2 delta (N1; note P1=1 ALONE does NOT bound ||P||, the +/- split does). [N1; GT-hyp; def-stochastic] \steptag{validated} \\[4pt]
\step{1.9} N9 (CONVERSE algebraic identity, exact ring computation). With Q from N8 and P the signed affine retraction (P\textasciicircum{}2=P EXACTLY, GT-hyp/def-stochastic), set D:=Q-P. Then Q=P+D, and in the associative unital ring M\_n(R) (N1 arena): Q\textasciicircum{}2=(P+D)\textasciicircum{}2=P\textasciicircum{}2+PD+DP+D\textasciicircum{}2=P+PD+DP+D\textasciicircum{}2 (using P\textasciicircum{}2=P). Subtract Q=P+D: Q\textasciicircum{}2-Q=(P+PD+DP+D\textasciicircum{}2)-(P+D)=PD+DP+D\textasciicircum{}2-D. Equivalently Q\textasciicircum{}2-Q=(Q-P)Q+P(Q-P)-(Q-P): expand (Q-P)Q+P(Q-P)=DQ+PD=D(P+D)+PD=DP+D\textasciicircum{}2+PD, minus (Q-P)=minus D, giving DP+D\textasciicircum{}2+PD-D=PD+DP+D\textasciicircum{}2-D, identical. This is an EXACT identity (no approximation), pure distributivity+associativity in M\_n(R), valid because P\textasciicircum{}2=P holds exactly. [N1; N8; P\textasciicircum{}2=P from GT-hyp/def-stochastic] \steptag{validated} \\[4pt]
\step{1.10} N10 (CONVERSE defect bound; closes CONVERSE). Take ||.|| (=||.|$|\_\{inf-\rangle$inf\}, N1) on the exact identity N9: Q\textasciicircum{}2-Q=PD+DP+D\textasciicircum{}2-D with D=Q-P. Triangle + submultiplicativity (N1): ||Q\textasciicircum{}2-Q||<=||P|| ||D||+||D|| ||P||... more precisely ||Q\textasciicircum{}2-Q||=||(Q-P)Q+P(Q-P)-(Q-P)||<=||Q-P|| ||Q|| + ||P|| ||Q-P|| + ||Q-P||. Substitute from N8: ||Q-P||=||P-Q||<=2 delta, ||Q||=1, ||P||<=1+2 delta. Hence ||Q\textasciicircum{}2-Q|| <= (2 delta)(1) + (1+2 delta)(2 delta) + 2 delta = 2 delta + 2 delta + 4 delta\textasciicircum{}2 + 2 delta = 6 delta + 4 delta\textasciicircum{}2. CONCLUSION (CONVERSE): row-normalising p\_i\textasciicircum{}+ yields a row-stochastic Q (N8) with ||P-Q||<=2 delta (N8) and ||Q\textasciicircum{}2-Q||<=6 delta+4 delta\textasciicircum{}2. [N1; N8; N9] \steptag{validated} \\[4pt]
\end{proofbox}

\end{document}
